\chapter{基础知识与方案依据}

身份基匹配加密并不是一种孤立的密码原语。它把分层身份基加密的访问控制能力与身份基签名的
来源认证能力组合起来，再由统一协议处理消息容量、对象编码和完整性保护。本章先介绍这两类
原语，再说明作品采用的格困难问题、具体算法和安全目标。第~3章在此基础上给出完整协议。

\section{从身份公钥到分层身份密钥}

传统公钥系统需要先取得用户公钥，再通过证书确认公钥属于谁。身份基密码把可读的身份字符串
直接作为公开标识，例如邮箱、设备编号或组织内账号。可信发钥机构持有主秘密，核验用户身份后
生成相应私钥。通信者只需知道对方身份和系统公开参数，不必为每个用户单独查询证书。

分层身份基加密（Hierarchical Identity-Based Encryption，HIBE）进一步把身份写成一条路径
$(\mathsf{id}_1,\ldots,\mathsf{id}_d)$。根机构可以为第一层身份生成密钥，上层用户还可以继续
为下一层身份委派密钥。在本作品使用的两层结构中，第一层放接收者$r$，第二层放发送者$s$：
\[
  (\mathsf{id}_1,\mathsf{id}_2)=(r,s).
\]
接收者取得一级密钥$\mathsf{dk}_r$后，可以为自己认可的发送者$s$生成叶密钥
$\mathsf{dk}_{r,s}$。加密者用公开参数和同一条身份路径产生密文。只有路径一致的叶密钥才能
恢复其中的短秘密。

一套两层HIBE包含五项基本操作：
\begin{enumerate}
  \item $\mathsf{Setup}$建立公开参数和主秘密；
  \item $\mathsf{Extract}$为第一层身份生成接收密钥；
  \item $\mathsf{Derive}$由接收密钥和第二层身份生成叶密钥；
  \item $\mathsf{Encrypt}$沿两层身份路径加密；
  \item $\mathsf{Decrypt}$使用叶密钥解密。
\end{enumerate}
这里最重要的性质是身份路径绑定。即使两名接收者使用同一套系统参数，他们的一级密钥也不能
互换；同一接收者为不同发送者得到的叶密钥同样不能互换。匿名HIBE还隐藏密文所对应的身份
路径，这正是Wang通用构造用于保护通信关系的基础~\cite{wang2024ibme}。

\section{身份基签名怎样认证发送者}

身份基签名（Identity-Based Signature，IBS）让验证者直接使用签名者身份和系统公开参数验签。
发钥机构首先建立IBS公开参数$\mathsf{ipk}$与主秘密$\mathsf{imsk}$，再为身份$s$生成签名密钥：
\[
 (\mathsf{ipk},\mathsf{imsk})\leftarrow\mathsf{Setup},
 \qquad
 \mathsf{sk}_s\leftarrow\mathsf{KeyExtract}(\mathsf{imsk},s).
\]
发送者用$\mathsf{sk}_s$签署消息$M$，接收者用$\mathsf{ipk}$、身份$s$和消息验签。实现还提供
签名加载方法，把传输字节恢复成严格类型的签名对象。

本作品接入了两种不同的IBS实现方式。第一种是原生格IBS，身份直接进入格上的目标、挑战或
陷门采样过程，Luo、SRGG、Sageloli和Wang--PW属于这一类。第二种是认证转换：发钥机构为用户
的普通签名公钥签发身份证书，用户再用对应私钥签消息；验证者依次检查机构证书和用户签名。
本作品按照BNN通用转换实现后一种方式，并接入ML-DSA与Falcon~\cite{galindo2006ibs,nist2024fips204,falcon2020}。

两种方式在公共协议中的职责相同，都是把发送身份$s$与待签字节绑定。区别主要体现在密钥生成、
签名尺寸和运行资源。认证转换可以直接利用成熟的普通签名实现；原生IBS则能保留论文给出的
身份结构和格关系。统一接口让这些差异停留在后端内部。

\section{四类格结构与困难问题}

格密码的安全性通常归结为在高维格中求解某类困难问题。本作品涉及的算法分布在普通矩阵格、
环格、模格和NTRU结构格上。它们共享线性代数、误差采样和短向量等工具，但对象形状与可利用的
代数结构不同。

\subsection{LWE与SIS的基本形式}

学习带误差问题（Learning With Errors，LWE）从带小误差的线性样本
\[
  (A,\;A\mathbf{s}+\mathbf{e}\bmod q)
\]
中隐藏秘密$\mathbf{s}$。误差$\mathbf{e}$使线性方程不能直接求解。格加密通常利用这一结构
隐藏明文或随机种子。短整数解问题（Short Integer Solution，SIS）则给出公开矩阵$A$，要求
找到非零短向量$\mathbf{z}$使
\[
  A\mathbf{z}=\mathbf{0}\pmod q,
  \qquad \lVert\mathbf{z}\rVert\leq\beta.
\]
格签名常把消息哈希成一个目标向量，再用陷门采样一个短原像；伪造者若能产生两组适当短原像，
便可导出SIS型短关系。

普通矩阵格直接保存$A$及向量。环格把向量组织为商环
$R_q=\mathbb Z_q[x]/(x^N+1)$中的多项式，循环卷积取代稠密矩阵乘法；Ring-LWE和Ring-SIS由此
获得紧凑表示和快速变换。模格再把环元素组成小秩向量或矩阵，在普通格的灵活性与环格的结构性
之间取折中；对应问题通常记为Module-LWE和Module-SIS~\cite{lpr2013ring,langlois2015module}。

NTRU结构也在多项式环上运算。短多项式关系既能构造陷门，也能支持快速签名和分层委派。LATTE
使用NTRU结构陷门与Ring-LWE密文路线；Falcon则在NTRU格上构造普通签名。作品在分类时保留
各论文的具体假设与模型，不把“环”“模”简单当作安全等级标签。

\subsection{困难问题、原语与系统职责的对应}

表~\ref{tab:assumption-map}把报告中的主要实现放在同一张图谱中。HIBE决定谁能够恢复随机种子，
IBS决定恢复后的消息是否来自指定发送者。两边可以采用不同维度、模数和格结构，组合协议只要求
它们满足统一的身份、容量和编码接口。

\begin{table}[H]
\centering\small
\caption{主要原语、格结构与系统职责}
\label{tab:assumption-map}
\begin{tabularx}{\textwidth}{@{}>{\raggedright\arraybackslash}p{.15\textwidth}>{\raggedright\arraybackslash}p{.21\textwidth}>{\raggedright\arraybackslash}p{.24\textwidth}>{\raggedright\arraybackslash}X@{}}
\toprule
\rowcolor{black!8} 原语 & 主要格结构或困难问题 & 在作品中的对象 & 在IB-ME中的职责 \\
\midrule
LATTE HIBE & NTRU结构陷门、Ring-LWE路线 & 两层接收密钥、叶密钥、32字节种子胶囊 & 按$(r,s)$路径封装种子 \\
Lai HIBE & 理想格陷门、PLWE路线 & 逐层追加的身份块与委派陷门 & 按$(r,s)$路径封装种子 \\
BNN＋ML-DSA & Module-LWE／Module-SIS及认证转换 & 机构证书、用户公钥与ML-DSA签名 & 认证发送身份与协议消息 \\
BNN＋Falcon & NTRU结构格签名及认证转换 & 机构证书、用户公钥与Falcon签名 & 认证发送身份与协议消息 \\
Luo & Module-LWE／Module-SIS & 模格公钥、响应和提示信息 & 原生身份签名 \\
SRGG & Module-LWE／Module-SIS & 身份目标、扰动承诺与短响应 & 原生身份签名 \\
Sageloli & Ring-SIS & 环陷门、身份盐、消息盐与短响应 & 原生身份签名 \\
Wang--PW & 矩阵SIS & G陷门、身份盐、消息盐与短原像 & 原生身份签名 \\
\bottomrule
\end{tabularx}
\end{table}

这张对应表说明了“多类格困难假设”在系统中的具体位置：它不是同时把多种假设堆在一个数学
公式中，而是让来自不同格结构的HIBE和IBS通过同一组合层工作。每套组合的密码性质仍由其
两项底层原语共同决定。

\section{通用构造的安全目标}

Wang等人的通用构造把两层匿名HIBE与IBS组合为IB-ME~\cite{wang2024ibme}。本作品采用其中的
普通隐私、加强认证与自我伪造防护目标。它们分别回答三个问题。

\textbf{消息和身份是否受到保护。} 未持有匹配接收密钥的参与者不能恢复消息；匿名HIBE还使
密文不直接暴露身份路径。第一层接收身份和第二层发送身份共同决定解密位置。

\textbf{发送来源是否可信。} 接收者恢复载荷后，以自己期望的发送身份验签。攻击者不能把一份
有效密文改写成另一发送者的消息，也不能只替换身份字段继续通过验证。

\textbf{合法用户能否冒充其他关系。} 一个拥有自己密钥的用户也不应借此构造另一组发送、接收
身份关系。HIBE路径与IBS签名消息同时包含双方身份，使“谁接收”和“认可谁发送”落在同一次
恢复过程中。

在软件中，这些目标被拆成三层实现。底层HIBE和IBS按照各自论文或标准运行；组合层固定身份
顺序、签名输入和解密次序；传输层再对完整密文字节做一次性认证。三层各自解决不同问题，
后两层不会修改底层格算法。

\section{已实现算法与代表组合}

日常组合矩阵由两种HIBE与五种IBS组成。两种HIBE是LATTE与Lai；五种IBS是BNN＋ML-DSA、
BNN＋Falcon、Luo模格、SRGG和Sageloli。十种配对均通过公共接口完成建立、发钥、加密、匹配
解密和身份错配处理。Wang--PW是随后完成的第六种IBS，其目标规模实现已经与LATTE完成一套
端到端组合；它的资源需求明显高于常规十组，因此单独列入第~4、6章。

报告选取A、B、C三组代表展开说明：
\begin{itemize}
  \item A组为LATTE＋BNN/ML-DSA-44，展示标准化普通签名经认证转换接入后的低延迟基线；
  \item B组为LATTE＋SRGG PARAM-I紧凑格式，展示原生模格IBS及对象编码优化；
  \item C组为Lai C2＋Sageloli PARAM-I，展示理想格HIBE、环格IBS和高成本委派场景。
\end{itemize}

三组不是简单的快慢排名。选择时同时考虑完整流程能否运行、格结构覆盖、在线时延、密钥和密文
尺寸、峰值内存、恢复能力以及能否体现作品中的关键优化。A适合高频消息基线，B用于观察原生
IBS和通信尺寸，C用于研究高成本原语怎样通过结构化计算与会话复用进入完整系统。Wang--PW则
补充矩阵SIS目标规模路线。

参数说明采用各原语自己的论文或标准口径。ML-DSA-44对应NIST安全类别2；SRGG PARAM-I论文
给出的量子密钥恢复与伪造成本分别约为95 bit和93 bit；Lai C2的现有系数实例量子估计约
154.2 bit；Wang--PW目标档按项目确定的约115 bit量子安全工作值使用。第~6章比较性能时保留
这些参数差异，不把不同原语的位数换算成一条统一排名。

\section{本章小结}

本章给出了作品的理论与算法坐标。两层HIBE沿$(r,s)$路径控制随机种子的恢复，IBS把发送身份
与协议消息绑定；普通矩阵格、环格、模格和NTRU结构格分别进入不同后端。统一组合不要求两项
原语共享参数，只要求它们遵守稳定接口。下一章将把这些原语组装成五个完整IB-ME算法，并说明
混合加密、完整密文认证和错误处理的具体位置。
